\abstract{\textbf{Purpose:} Estimated Time of Arrival (ETA) prediction underlies route planning, dispatch, and arrival notifications, yet many systems estimate ETA from traffic history without explicit knowledge of a vehicle's planned route. This work examines whether a vehicle-centered dynamic graph with route awareness improves ETA prediction.

\textbf{Methods:} We present DSTRA-GNN (Dynamic Spatio-Temporal Route-Aware GNN), which integrates vehicle-level dynamics and route intent in a Mixture-of-Experts framework. Vehicles are graph nodes, interaction edges model local traffic context, and a GRU-based temporal aggregator processes sequential graph snapshots. We evaluate on SUMO simulation data and Porto-G, a graph representation of the Porto taxi trajectory dataset (ECML-PKDD 2015).

\textbf{Results:} Ablations over six variants show consistent, compounding gains from route awareness and temporal context. On duration-matched Porto-G, MAE drops from 93.3\,s (static graph only) to 68.1\,s (27\% reduction), performing competitively with trajectory-based baselines such as DeepTTE (68.6\,s) and MetaTTE (69.5\,s). On SUMO, MAE drops from 80.1\,s to 54.8\,s (32\% reduction) under the shared protocol range; on the broader p95-filtered SUMO population, the full model remains the strongest evaluated variant and outperforms applicable classical and graph-native baselines at 72.4\,s MAE.

\textbf{Conclusion:} Vehicle-centered dynamic graphs with explicit route state improve ETA prediction across simulated and real-world data, positioning graph-snapshot ETA modeling as a competitive alternative to trajectory-native sequence models.}
